\begin{claim}{F.22}{Proved}{The actual two-point kernel, its compact generators, and the remaining source map}
Retain the original F.19 base configuration. Let $p=a(s_0)$,
$n=-a(s_0)$, and let $S'$ be the plane with its seven old
finite labels removed. Filling the four labels
$U=\{v_u:u\ne\zeta\}$ gives the plane $S$ with the three
old finite labels $0,1,v_\zeta$ removed. The actual
two-point groups are
\[
 J=\widehat{\pi_1(\operatorname{Conf}_2(S'),(p,n))},\qquad
 J_4=\widehat{\pi_1(\operatorname{Conf}_2(S),(p,n))},\qquad
 D^+=\ker(J\to J_4).
                                                               \tag{F.22.1}
\]
There are compatible based point-pushing sections of the
two positive-point projections, described below, and
geometric free coordinates in which
\[
 \begin{gathered}
 V_+=\widehat F(g_1,g_2,g_3,u_1,u_2,u_3,u_4),\qquad
 Q_+=\widehat F(g_1,g_2,g_3),\\
 V=\widehat F(X_1,X_2,X_3,U_1,U_2,U_3,U_4,t),\qquad
 P=\widehat F(X_1,X_2,X_3,t),\\
 R_+=\ker(V_+\to Q_+),\qquad R_1=\ker(V\to P).
 \end{gathered}                                                \tag{F.22.2}
\]
The displayed fillings kill precisely $u_j$ and $U_j$,
respectively. With the actual section $\sigma:V_+\to J$
and fiber inclusion $\iota:V\to J$, there are mutually
inverse continuous maps
\[
 \begin{aligned}
 R_1\rtimes R_+&\xrightarrow{\ \Theta\ }D^+,
       &(b,g)&\longmapsto\iota(b)\sigma(g),\\
 D^+&\longrightarrow R_1\rtimes R_+,
       &d&\longmapsto
       \bigl(\iota^{-1}(d\sigma(p_+d)^{-1}),p_+d\bigr).
 \end{aligned}                                                \tag{F.22.3}
\]
Here $p_+$ forgets the negative point; the action in this
semidirect product is the actual conjugation by $\sigma$.
Its complete free-word recipe is (F.22.11)--(F.22.12).

For a profinite space $T$, write $\mathfrak F(T)$ for the
free profinite group on continuous maps from $T$, as
constructed in the proof. This definition retains the
topology on the generator space. Put $I=\{1,2,3,4\}$.
Then the exact compact-generator identifications are
\[
 \begin{aligned}
 \mathfrak F(Q_+\times I)&\xrightarrow{\ \sim\ }R_+,
       &(q,j)&\longmapsto s_+(q)u_js_+(q)^{-1},\\
 \mathfrak F(P\times I)&\xrightarrow{\ \sim\ }R_1,
       &(v,j)&\longmapsto s_1(v)U_js_1(v)^{-1},
 \end{aligned}                                                \tag{F.22.4}
\]
where $s_+$ and $s_1$ retain the displayed free generators.

The other one-point forgetting is the actual
$p_-:J\to V_-$, with
\[
 \begin{gathered}
 V_-=\widehat F(X_1,X_2,X_3,U_1,U_2,U_3,U_4),\qquad
 Q_-=\widehat F(X_1,X_2,X_3),\\
 R_-=\ker(V_-\to Q_-),\qquad
 W=D^+\cap\ker p_+\cap\ker p_- .
 \end{gathered}                                                \tag{F.22.5}
\]
One has the exact sequence
\[
 1\longrightarrow W\longrightarrow D^+
        \xrightarrow{(p_+,p_-)}R_+\times R_-\longrightarrow1.
                                                               \tag{F.22.6}
\]
Inside the actual fiber $V$, let $s_-$ retain $Q_-$ and put
$\eta(q,e)=s_-(q)e s_-(q)^{-1}$ for
$e\in\{U_1,U_2,U_3,U_4,t\}$. Then
\[
 W=
 \overline{\langle\!\langle
       [\eta(q,U_j),\eta(q',t)]
       :q,q'\in Q_-,\ j\in I
                         \rangle\!\rangle}_{V} .
                                                               \tag{F.22.7}
\]
The literal element $[U_1,t]$ belongs to $W$ and
generates a full closed copy of $\widehat{\mathbb Z}$.

Finally let $\Delta$ be the sixteen actual source
collisions, with the F.19 source meridians and section
$j:L\to H^\dagger$. There is a specified continuous map
\[
 \Lambda:\mathfrak F(L\times\Delta)\longrightarrow D^+,
       \qquad(l,d)\longmapsto j(l)\nu_dj(l)^{-1},
       \qquad\operatorname{im}\Lambda=K^+ .
                                                               \tag{F.22.8}
\]
Its coordinates in (F.22.3) are given by the actual
two projection formulas there, applied to the displayed
element. Thus the remaining equality $K^+=D^+$ is
exactly surjectivity of this named continuous homomorphism
onto the displayed semidirect product. This statement
does not settle that surjectivity or the original
residual source-membership equation.
Write $\Delta_+$ and $\Delta_-$ for the eight collisions
with the positive and negative new point, respectively,
and put $\mathcal K_\pm=p_\pm(K^+)$. The exact joint
source image is their full direct product:
\[
 1\longrightarrow K^+\cap W\longrightarrow K^+
       \xrightarrow{(p_+,p_-)}\mathcal K_+\times\mathcal K_-
       \longrightarrow1 .
                                                              \tag{F.22.8a}
\]
Each marginal is the actual image of the corresponding
eight-collision source filling kernel, as specified
in (F.22.21). Their equality with $R_+$ and $R_-$
remains unresolved.
\end{claim}

\begin{proof}
\emph{Compatible sections at the original basepoint.}
All seven old finite points, $p$, and $n$ are distinct
by F.17. Choose an embedded finite tree in the plane
joining the seven old points and $p$ while avoiding
$n$. A sufficiently small regular neighborhood is a
closed disk $B$ containing those eight points in its
interior and excluding $n$. The inclusions
\[
 B^\circ\setminus\{\hbox{seven old points}\}\longrightarrow S',
 \qquad
 B^\circ\setminus\{\hbox{three retained old points}\}
                    \longrightarrow S
                                                               \tag{F.22.9}
\]
induce isomorphisms on fundamental groups based at
the original point $p$. Indeed a collar of the disk
boundary and the outside annulus retract onto that
boundary, so adjoining the outside region changes
none of the punctured-disk fundamental-group generators
or relations.

On each punctured disk the map $z\mapsto(z,n)$ is an
actual based map to its ordered two-point configuration
space. Compose its fundamental-group map with the
inverse of the relevant isomorphism (F.22.9).
This gives sections $\sigma_{\mathrm d}$ for $S'$ and
$\sigma_{4,\mathrm d}$ for $S$.
The two disk inclusions and the two maps $z\mapsto(z,n)$
commute with filling the four old punctures. Their
inverse group isomorphisms therefore commute too.
Hence these very sections form a commuting square
with the actual filling $J_{\mathrm d}\to J_{4,\mathrm d}$.
All basepoints here are $(p,n)$; the comparison paths
to the F.19 two-point coordinates are the constant
paths at that configuration.

The Fadell--Neuwirth sequences have kernels
$\pi_1(S'\setminus\{p\},n)=F_8$ and
$\pi_1(S\setminus\{p\},n)=F_4$.
Their injections follow from $\pi_2(S')=\pi_2(S)=0$.
The full-topology and completed-kernel comparisons
are exactly those established in F.19 for these
two-point configuration groups. The actual discrete
sections extend to full completions, their section
identities persist, and their filling square still
commutes. Thus
\[
 J=V\rtimes_\sigma V_+,\qquad
 J_4=P\rtimes_{\sigma_4}Q_+,
 \qquad
 (b,g)\longmapsto(\overline b,\overline g)
                                                               \tag{F.22.10}
\]
is the actual filling map in these coordinates.
The equation $\overline g=1$ puts $g$ in $R_+$,
and then $\overline b=1$ puts $b$ in $R_1$.
Multiplication and the section inverse give exactly
the two maps in (F.22.3).

\emph{Free coordinates and the retained Artin action.}
Choose a geometric chart on $B$ in which the old
retained points occupy positions $1,2,3$, the four
other old points occupy positions $4,5,6,7$, and $p$
occupies position $8$. The exterior point $n$ is
the basepoint for puncture meridians. Such a chart
is an orientation-preserving homeomorphism of the
disk with the stated finite labelled set: move
the points along polygonal arcs avoiding the other
points, using successive small disk isotopies.
The induced fundamental-group isomorphisms have
the inverse chart as their literal inverse.
Both filling maps use this same chart and its
restriction to the retained labelled set.
They therefore commute with the actual coordinate
maps, including the basepoints.

Write $X_4=U_1,\ldots,X_7=U_4,X_8=t$.
Use the literal clockwise exchanges and the Artin
substitutions calibrated in F.20.7--9 on
$\widehat F(X_1,\ldots,X_8)$.
For the positive point moving once clockwise around old
position $i$ along its lower-half-plane whisker, the
based geometric free generator $g_i$ has the braid
\[
 \beta_i=\sigma_7\cdots\sigma_{i+1}\sigma_i^2
                  \sigma_{i+1}^{-1}\cdots\sigma_7^{-1},
          \qquad1\leq i\leq7,
 \qquad u_j=g_{j+3}.
                                                               \tag{F.22.11}
\]
The descending word carries the last strand below
the intervening points, its square is the clockwise
lasso, and the inverse word returns the same whisker.
Deleting that strand gives the trivial old braid by
contracting the old motion in its projection-order
chambers. These seven lower-whiskered lassos form the
actual geometric kernel basis. The section keeps the
negative point $n$ fixed outside $B$, so the exact
fibre-coordinate comparison below makes its conjugation
on the negative-point fiber
\[
 \alpha_{g_i}=\phi_{\beta_i},
 \qquad
 \phi_{\sigma_k}(X_k)=X_kX_{k+1}X_k^{-1},
 \quad\phi_{\sigma_k}(X_{k+1})=X_k,
 \quad\phi_{\sigma_k}(X_j)=X_j\ (j\ne k,k+1).
                                                               \tag{F.22.12}
\]
The clockwise braid word is composed in the same
chronological and substitution order as F.20, and its
inverse word gives its inverse map.

Here is the precise point-pushing identification
underlying that conjugation formula. For a pure
braid on the eight punctures and the exterior ninth
point, let $d$ be the ninth Artin meridian and write
$\phi_b(d)=W_bdW_b^{-1}$. Erasing $d$ gives $w_b$;
put $u_b=w_b^{-1}$. Exactly as in F.20.8, substitution
gives $u_{bc}=u_b\phi_b(u_c)$ after erasing the ninth
strand. The section on the first eight strands has
$u=1$. Write $\gamma_i$ for the eight lower-whiskered
clockwise point-pushing lassos of the ninth point,
given by the same braid pattern with $\sigma_8$ as
its first exchange. Their values of $u$ are
\[
       \tau_i=X_8^{-1}\cdots X_{i+1}^{-1}
                    X_i^{-1}X_{i+1}\cdots X_8,
                \qquad1\leq i\leq8.
                                                               \tag{F.22.13}
\]
This triangular substitution is its own inverse:
starting with $i=8$ and substituting downward,
each adjacent inverse pair cancels and the image
of $\tau_i$ is $X_i$. The actual forward coordinate
map is $\gamma_i\mapsto\tau_i(X)$; its inverse is
$X_i\mapsto\tau_i(\gamma)$. Both composites are
the identity on the named free bases. They therefore
give the actual free-group isomorphism and its full
profinite inverse between the fiber and these Artin
coordinates. The crossed multiplication formula now
makes conjugation by the actual section exactly
(F.22.12). Erasing any of the four old labels,
or erasing $t$, commutes with (F.22.13); the
remaining triangular formula has exactly the
remaining labels in their retained order.
This proves the coordinate form of both fillings.

The plane braid groups in this calculation map
to the original normalized moduli groups through
the explicit coordinate isomorphism
\[
 (z_0,z_1,z_2,\ldots)
   \longmapsto
 \left(z_0,z_1-z_0,
           \left(\frac{z_j-z_0}{z_1-z_0}\right)_{j\geq2}\right).
                                                               \tag{F.22.14}
\]
Its inverse multiplies the last coordinates by
$z_1-z_0$ and adds $z_0$. The first two old labels
are fixed on every loop in $J$, so their two
affine coordinates are constant. This is the
actual comparison with the normalized F.19
configuration groups.

For completeness, the action extends continuously
from $F_7$ to $V_+=\widehat F_7$ on $V=\widehat F_8$.
Every open subgroup of this finitely generated
profinite $V$ contains a characteristic open
subgroup: intersect its finitely many subgroups
of at most the given index. The action on each
resulting finite characteristic quotient has
finite image and extends through $V_+$.
The compatible inverse automorphisms give the
continuous inverse action. This is also the
action already supplied by the actual section
in (F.22.10).

\emph{A free profinite group on a compact generator space.}
For a profinite space $T$, begin with the abstract
free group on its underlying set. Take its inverse
limit over finite quotients for which the map
from $T$ to that finite quotient is continuous;
denote the result by $\mathfrak F(T)$.
Products of finitely many such quotient maps again
have continuous generator maps, so these quotients
form a directed system. The canonical map
$e:T\to\mathfrak F(T)$ is continuous and its
image generates a dense subgroup. Every continuous
map from $T$ to any profinite group $H$ extends
uniquely to a continuous homomorphism
$\mathfrak F(T)\to H$: compose with every finite
quotient of $H$, extend on the abstract free
generators, and take the compatible inverse limit.
This construction proves the precise universal
property used here.

Let $Q$ be one of the finitely generated free
profinite groups in (F.22.2), and let $I$ be the
finite set of the additional generators.
Left multiplication on $Q\times I$ induces
automorphisms of $\mathfrak F(Q\times I)$.
This action is jointly continuous. Indeed for
any continuous map $Q\times I\to A$ to a finite
group, compactness gives an open normal subgroup
$N\triangleleft Q$ through which all its finitely
many coordinate functions factor. Its translates
therefore form a finite set indexed by $Q/N$.
Intersect the kernels of the corresponding
finite quotient homomorphisms of $\mathfrak F$.
The intersection is open and $Q$-invariant, and
the action on this quotient factors through a
finite quotient of $Q$. These invariant open
subgroups refine every finite quotient kernel.
They prove joint continuity and give the actual
profinite semidirect product
$\mathfrak F(Q\times I)\rtimes Q$.

Let $\widehat F(Q;I)$ denote the free profinite
group obtained by adjoining these finitely many
extra generators $v_j$ to the fixed free basis
of $Q$. Its retraction $\rho$ kills the $v_j$,
and its section $s$ retains the basis of $Q$.
There are inverse continuous homomorphisms
\[
 \begin{aligned}
 \widehat F(Q;I)&\longrightarrow
                     \mathfrak F(Q\times I)\rtimes Q,\\
 x&\longmapsto(1,x)\quad(x\text{ in the retained basis}),
       \qquad v_j\longmapsto(e(1,j),1),\\
 \mathfrak F(Q\times I)\rtimes Q
       &\longrightarrow\widehat F(Q;I),\\
 e(q,j)&\longmapsto s(q)v_js(q)^{-1},
       \qquad q\longmapsto s(q).
 \end{aligned}                                                \tag{F.22.15}
\]
The second generator map is continuous and
$Q$-equivariant, so the universal property and
the semidirect-product law give its homomorphism.
Both composites fix every retained generator,
every $v_j$, and every $e(q,j)$ by the displayed
conjugation formula. These sets generate densely,
so both composites are the identity. The
projections onto $Q$ are precisely $\rho$ and
the semidirect-product projection. Their kernels
are therefore identified by these inverse maps.
Apply this construction first to $(Q_+,u_j)$
and then to $(P,U_j)$ to prove (F.22.4).
No full completion of an infinitely generated
discrete kernel has been substituted in this
argument.

The action in (F.22.3) can now be evaluated
on every compact generator, not only on a
finite selection. The element $e(q,j)$ of
$R_+$ has image $s_+(q)u_js_+(q)^{-1}$ in
$V_+$. Its action is
$\alpha_{s_+(q)}\alpha_{u_j}\alpha_{s_+(q)}^{-1}$.
For a compact generator $e(v,k)$ of $R_1$,
apply that automorphism to
$s_1(v)U_ks_1(v)^{-1}$ in $V$, then use the
inverse of (F.22.15). The result lies in
$R_1$ because the induced action of $R_+$
on $P$ is the identity. The formula is
continuous in both $q$ and $v$ by the
actual continuous group operations and the
proved action continuity.

\emph{The interaction kernel retained by both one-point forgettings.}
On the fiber $V$, the negative-point projection
$p_-$ fills the positive puncture, so its map
$h:V\to V_-$ kills $t$ and retains every other
letter. On the section $\sigma(V_+)$ it is
the identity element: that section keeps the
negative point fixed. Hence
\[
        (p_+,p_-)(\Theta(b,g))=(g,h(b)).
                                                               \tag{F.22.16}
\]
The restriction $h:R_1\to R_-$ is onto.
For an element of $R_-$, retain its seven
letters in $V$; its old-puncture filling is
the identity, so it belongs to $R_1$ and
maps back to that element under $h$.
Together with the arbitrary $g\in R_+$,
this proves surjectivity in (F.22.6), with
kernel $W=R_1\cap\ker h$ inside $V$.

Apply (F.22.15) with retained group $Q_-$
and the five extra letters $U_1,\ldots,U_4,t$.
The kernel of $V\to Q_-$ is
\[
 \mathcal R=
 \mathfrak F\bigl((Q_-\times I)\sqcup(Q_-\times\{t\})\bigr)
       =\mathcal A\amalg\mathcal B,
                                                               \tag{F.22.17}
\]
where $\mathcal A=\mathfrak F(Q_-\times I)$,
$\mathcal B=\mathfrak F(Q_-\times\{t\})$,
and the coproduct equality follows directly
from the continuous-map universal property
on the two compact open-and-closed pieces.
The two fillings restrict to the retractions
$\mathcal R\to\mathcal B$ and
$\mathcal R\to\mathcal A$. Consequently
\[
             W=\ker(\mathcal A\amalg\mathcal B
                            \longrightarrow\mathcal A\times\mathcal B).
                                                               \tag{F.22.18}
\]
Quotienting $\mathcal R$ by the closed normal
closure of all the displayed generator
commutators makes its two factor images
commute. The multiplication map from
$\mathcal A\times\mathcal B$ to this quotient
is then a homomorphism. Its composite with
the two retractions is the identity, and
the reverse composite fixes all generators.
These are inverse homomorphisms. Thus their
kernel is exactly the stated normal closure.
It is invariant under $Q_-$, whose action
sends $(q,q')$ to $(aq,aq')$.
Its normal closure in $V=\mathcal R\rtimes Q_-$
is therefore the same subgroup. This proves
(F.22.7), with its actual ambient group.

For every integer $m\geq2$, send $U_1$ to
$I+E_{12}$, send $t$ to $I+E_{23}$, and send
the other free generators of $V$ to the
identity in the finite group
$\operatorname{UT}_3(\mathbb Z/m)$.
These assignments define a continuous
homomorphism. With $[a,b]=aba^{-1}b^{-1}$,
\[
       [U_1,t]\longmapsto I+E_{13},\qquad
       [U_1,t]^z\longmapsto I+(z\bmod m)E_{13}
                \quad(z\in\widehat{\mathbb Z}).
                                                               \tag{F.22.19}
\]
The commutator is killed by both fillings,
so it belongs to $W$. The displayed finite
images detect every nonzero
$z\in\widehat{\mathbb Z}$ through some modulus.
The exponent map to its closed cyclic
subgroup is therefore injective as well
as onto, proving the asserted full
procyclic subgroup.

\emph{The actual compact-generator source image.}
The F.19 source section retains its eleven
old free meridians in the rank-twenty-seven
source; the sixteen additional free meridians
are $m_d$, $d\in\Delta$. Apply (F.22.15)
to that very based free system. It identifies
\[
 \ker(q:L^\dagger\to L)
       \cong\mathfrak F(L\times\Delta),
       \qquad(l,d)\longmapsto i(l)m_di(l)^{-1}.
                                                               \tag{F.22.20}
\]
Composition with the actual map
$\Phi^\dagger:L^\dagger\to G$ gives exactly
(F.22.8), because $\Phi^\dagger i=j$ and
$\Phi^\dagger(m_d)=\nu_d$ with the retained
source whisker. F.19.10 identifies the image
of this source kernel with $K^+$; alternatively
its image is precisely the kernel of the
induced onto map $H^\dagger\to L$, since
the commuting source filling has that
kernel. The map in (F.22.8) is continuous
and its image is closed, since its domain
is compact.

\emph{The two marginal source images and their joint image.}
Put $Y^\pm=Y\setminus\Delta_\pm$ and
$L_\pm=\widehat{\pi_1(Y^\pm,s_0)}\cong\widehat F_{19}$.
The source fillings $q_\pm:L_\pm\to L$ retain
the eleven old basis elements and kill their
eight added meridians.
Let $F_+:G\to P_9^+$ forget the negative point,
and let $F_-:G\to P_9^-$ forget the positive point.
These are the actual full nine-point forgetting
maps; their restrictions to $J$ are $p_+$ and
$p_-$ with the displayed one-point kernel targets.
The rational source configuration after forgetting
the negative point is defined on all of $Y^+$;
the only extra singularities of that map are
the eight positive collisions. The corresponding
statement holds with the signs interchanged.
Thus the actual source maps factor as
\[
 \begin{gathered}
 F_\pm\Phi^\dagger
       =\Phi_\pm r_\pm,
       \qquad r_\pm:L^\dagger\twoheadrightarrow L_\pm,\\
 \mathcal K_\pm
       =\Phi_\pm(\ker q_\pm)
       =\operatorname{im}
        \bigl(\mathfrak F(L\times\Delta_\pm)
                                \longrightarrow R_\pm\bigr).
 \end{gathered}                                                \tag{F.22.21}
\]
Here $r_+$ fills the negative collision points,
and $r_-$ fills the positive ones.
These equalities retain their based source maps:
the signed free basis embeds $L_\pm$ back into
$L^\dagger$ by its named old and signed meridians,
so $r_\pm(\ker q)=\ker q_\pm$.
The last map sends $(l,d)$ to
$F_\pm(j(l)\nu_dj(l)^{-1})$.

For $d\in\Delta_+$, forgetting the positive
point removes the collision. The remaining
configuration map extends across $d$, so the
based source meridian has $F_-(\nu_d)=1$.
For $d\in\Delta_-$ one similarly has
$F_+(\nu_d)=1$.
For every $l\in L$ these identities give
\[
 \begin{aligned}
 F_-(j(l)\nu_dj(l)^{-1})&=1\quad(d\in\Delta_+),\\
 F_+(j(l)\nu_dj(l)^{-1})&=1\quad(d\in\Delta_-).
 \end{aligned}                                                \tag{F.22.22}
\]
The maps here are defined on all of $G$;
the calculation does not apply $p_\pm$ to
an element outside $J$.

Let $K_+^{\mathrm{src}}$ and $K_-^{\mathrm{src}}$
be the compact images of
$\mathfrak F(L\times\Delta_+)$ and
$\mathfrak F(L\times\Delta_-)$ in $K^+$
under (F.22.8). Their joint images are
$\mathcal K_+\times\{1\}$ and
$\{1\}\times\mathcal K_-$, respectively,
by (F.22.21)--(F.22.22).
Given a pair in $\mathcal K_+\times\mathcal K_-$,
lift its first coordinate to $K_+^{\mathrm{src}}$
and its second to $K_-^{\mathrm{src}}$.
Their product in $K^+$ maps to the given pair.
This proves (F.22.8a). Its kernel is
$K^+\cap W$ by (F.22.6), giving the exact
commuting diagram of actual inclusions
\[
 \begin{array}{ccccccccc}
 1&\longrightarrow&K^+\cap W&\longrightarrow&K^+
     &\longrightarrow&\mathcal K_+\times\mathcal K_-
     &\longrightarrow&1\\
 &&\downarrow&&\downarrow&&\downarrow\\
 1&\longrightarrow&W&\longrightarrow&D^+
     &\longrightarrow&R_+\times R_-
     &\longrightarrow&1 .
 \end{array}                                                  \tag{F.22.23}
\]
The first exact row concerns the proved marginal
images, with no enlargement of either marginal.

The point-pushing section $\sigma$ and the
source section $j$ have both been retained
as their actual maps. To find the coordinates
of any generator in (F.22.8), set
$z=j(l)\nu_dj(l)^{-1}$, compute the literal
one-point forgetting $g=p_+(z)$, and take
$b=\iota^{-1}(z\sigma(g)^{-1})$.
This is exactly $\Theta^{-1}(z)$, with no
identification of the two sections.
The first coordinate belongs to $R_1$ and
the second to $R_+$ by the proved inverse
maps. These formulas identify the remaining
source-membership problem in the displayed
kernel presentation.

The exact producer
\path{code/part_a/joint_kernel_presentation_f22.py}
checks all seven Artin action words and
their inverses, their pure conjugators,
their boundary product, and their two
actual filling squares. It also records
the deleted-label conjugators after the
old filling, checks that $[U_1,t]$ is
killed by both fillings, and checks the
central power formula in thirty finite
Heisenberg groups. The algebraic formula
(F.22.19), proved for every modulus,
supplies the full profinite conclusion.
The compact-generator construction and
all semidirect-product inverse maps are
proved above independently of those
finite checks.
\end{proof}
